\section{Development of the V2 Prototype}
\label{sec:intro-v2}

The first prototype was used to develop the power circuit and burst
controller. V2 was then designed with the MCU on the converter board in place of
the separate controller platform. The charge-pump
PFC and isolated resonant paths were retained, together with the supplied
planar magnetics and downstream USB-C output stage.

The main V2 design specifications are summarized in
Table~\ref{tab:intro-v2-specs}. The circuit connections and controller
timing are described in \cref{sec:hw-v2}; the complete schematic and PCB
print set are included in \cref{sec:app-v2-oct27}.

\begin{table}[htbp]
\centering
\caption{V2 design specifications.}
\label{tab:intro-v2-specs}
\small
\begin{tabularx}{\textwidth}{@{}p{0.25\textwidth}X@{}}
\toprule
Function & V2 implementation\\
\midrule
Controller & STM32F410CBU6 integrated on the converter PCB\\
Power stage & LMG2652RFBR GaN half bridge shared by the charge-pump and resonant paths\\
Feedback & Scaled primary-voltage input (PA0), isolated secondary-bulk voltage feedback through AMC0330 (PA1), and primary half-bridge-current sensing (PA2)\\
USB-C output & MPQ4242 output stage downstream of the isolated secondary bulk bus\\
Local supplies & BAC1S12SC bias supply and separate LMR54410 regulators for primary and secondary local supplies\\
PCB & Four copper layers, including internal GND and PWR layers\\
Carrier timing & Default \SI{1}{\mega\hertz} carrier and \SI{40}{\nano\second} programmed dead time; both adjustable with output disabled (\cref{sec:ctrl-v2-adjustable-timing})\\
\bottomrule
\end{tabularx}
\end{table}

On V2, carrier commands use PA8/PB13, burst enable uses PA11, and UART
TX/RX use PA9/PA10. The MCU continues to regulate the secondary bulk supply,
which feeds the USB-C PD output stage.

During later powered V2 tests, heating was observed in the external
series resonant inductor. It was removed, and the tank was changed to use
transformer leakage with approximately \SI{1.66}{\nano\farad} resonant
capacitance. The software dead-time setting was approximately
\SI{60}{\nano\second}. The input bridge continued to heat up during the
test labelled \SI{95}{\watt} (\cref{sec:res-v2-testing}).

The V2 drawings show the resonant branch before the external inductor was
removed. Some component values and feedback scales still need checking
against the assembled board (\cref{sec:hw-v2}). Further tests are needed
to resolve the bridge heating and compare the fitted V2 converter with
the V3 edge-card design.
